\chapter{The \texorpdfstring{$p$}{p}-adic Numbers}\label{ch:padic}

\begin{center}\itshape
In the $2$-adic world, $1+2+4+8+\cdots=-1$, and $-7$ has a square root.\\
Both statements are true, and they say that ``being close'' can mean ``being congruent modulo a high power of $p$''.
\end{center}

In 1897 Kurt Hensel had an idea that looked like a joke: a polynomial can be expanded in powers of $x-a$ around a point, so why not expand an integer in
powers of a prime $p$, and allow the expansion to go on forever? The result, the $p$-adic numbers, turned out to be no joke at all: they are now as basic
to number theory as the real numbers, and the rule ``think about each prime separately, then glue'' runs through all of modern arithmetic.

Fix a prime $p$. Volume~2 taught us to solve equations modulo $p$, then modulo $p^2,p^3,\dots$ by Hensel's lemma. The $p$-adic numbers are what you get when you never stop: a $5$-adic integer is a coherent list of residues modulo $5,25,125,\dots$, just as a real number is a coherent list of decimal approximations. In this chapter we build them, compute with them, and see why there are no other ways to measure size on $\Q$.

\begin{reminder}[How to read an infinite digit string]
We list $p$-adic digits starting at $p^0$, then $p^1$, then $p^2$.
Thus ``the first $n$ digits'' means a remainder modulo $p^n$, not an
approximation of small ordinary real error. Knowing more digits means knowing
more divisibility. The same sequence of rational numbers can converge in one metric and
fail to converge in another; every claim of convergence needs its metric.
\end{reminder}

\section{The \texorpdfstring{$p$}{p}-adic absolute value}

% t1-pass2 motive 11a826ac
Divisibility by $p$ can be read as closeness: two integers are close when their difference is divisible by a high power of $p$. The resulting absolute value is ultrametric, so geometry changes, and the product formula says that the real absolute value and all the $p$-adic ones are not independent. Chapter~9 will build the local fields on this absolute value, and Chapter~12 will build Tate's thesis on it.

For $x\in\Q^*$, write $x=p^v\frac ab$ with $p\nmid ab$; the integer $v=v_p(x)$ is the \emph{$p$-adic valuation} of $x$ (for integers, Volume~2). It is well defined, since $v_p(x)=v_p(\text{numerator})-v_p(\text{denominator})$ for any fraction representing $x$. The \emph{$p$-adic absolute value} is
\[
|x|_p=p^{-v_p(x)}\quad(x\ne0),\qquad|0|_p=0.
\]
So $x$ is $p$-adically small when it is divisible by a high power of $p$.

\begin{think}
Which of $1$, $25$, $\frac15$, $5^{100}$, $10^{100}$ is closest to $0$ for $|\cdot|_5$? Which two of the numbers $1$, $26$, $126$ are closest to each other?
\end{think}

\begin{proposition}\label{prop:ultra}
For $x,y\in\Q$:
\begin{enumerate}[label=(\alph*)]
\item $|xy|_p=|x|_p\,|y|_p$;
\item (ultrametric inequality) $|x+y|_p\le\max(|x|_p,|y|_p)$, with equality when $|x|_p\ne|y|_p$;
\item (product formula) for $x\ne0$, $\displaystyle|x|_\infty\prod_p|x|_p=1$, where $|x|_\infty$ is the usual absolute value and the product (over all primes) has finitely many factors different from $1$.
\end{enumerate}
\end{proposition}
\begin{proof}
(a) $v_p(xy)=v_p(x)+v_p(y)$, by uniqueness of factorisation of numerators and denominators. (b) We may assume $x,y,x+y\ne0$ and $v_p(x)\le v_p(y)$. Write $x=p^{v}\frac ab$ and $y=p^{v}\frac cd$ with $v=v_p(x)$, $p\nmid abd$ (and possibly $p\mid c$). Then $x+y=p^v\frac{ad+bc}{bd}$ with $p\nmid bd$, so $v_p(x+y)\ge v$, which is the inequality. If $v_p(y)>v$, then $p\mid c$, so $p\nmid ad+bc$ and $v_p(x+y)=v$. (c) Write $x=\pm\prod_pp^{v_p(x)}$; then $|x|_\infty=\prod_pp^{v_p(x)}$ and $\prod_p|x|_p=\prod_pp^{-v_p(x)}$.
\end{proof}

The ultrametric inequality makes $p$-adic geometry strange and pleasant. Define the distance $d(x,y)=|x-y|_p$ and the ball $B(a,r)=\{x:|x-a|_p\le r\}$.

\begin{proposition}\label{prop:balls}
\begin{enumerate}[label=(\alph*)]
\item Every triangle is isosceles: among $|x-y|_p$, $|y-z|_p$, $|x-z|_p$, the two largest are equal.
\item Every point of a ball is a centre: if $b\in B(a,r)$ then $B(b,r)=B(a,r)$.
\item Two balls are either disjoint or one contains the other.
\end{enumerate}
\end{proposition}
\begin{proof}
(a) $x-z=(x-y)+(y-z)$; by Proposition~\ref{prop:ultra}(b), if $|x-y|_p\ne|y-z|_p$ then $|x-z|_p$ equals the larger of the two. (b) If $|x-a|_p\le r$ then $|x-b|_p\le\max(|x-a|_p,|a-b|_p)\le r$; the converse is symmetric. (c) If $B(a,r)$ and $B(b,s)$ with $r\le s$ share a point $c$, then by (b) they are $B(c,r)\subseteq B(c,s)$.
\end{proof}

In $\Z$, the ball $B(a,p^{-k})$ is the residue class of $a$ modulo $p^k$. So the balls form a tree: each class modulo $p^k$ splits into $p$ classes modulo $p^{k+1}$ (Figure~\ref{fig:tree}).

\begin{figure}[ht]
\centering
\fig{ch12-tree}
\caption{The $2$-adic integers as an infinite binary tree. A node at depth $k$ is a residue class modulo $2^k$ (a ball of radius $2^{-k}$), and an infinite path from the root is a $2$-adic integer, read off as its binary digits. The dashed path, all digits $1$, is $-1$; the solid path that branches from it is $\frac13=1+2+0\cdot4+8+0\cdot16+\cdots$.}
\label{fig:tree}
\end{figure}

\begin{figure}[ht]
\centering
\fig{ch12-ruler}
\caption{$|n|_2$ for $1\le n\le 64$: the ``ruler function''. Most integers have $2$-adic size $1$; the multiples of high powers of $2$ are the small ones.}
\label{fig:ruler}
\end{figure}

\subsection*{A different ruler, with numerical readings}
For $x=75/14$ we have $v_5(x)=2$, so $|x|_5=1/25$, while
$v_7(x)=-1$, so $|x|_7=7$. Ordinary size, $5$-adic size and $7$-adic size
answer different questions about the same number.
Also
\[
|26-1|_5=5^{-2},\qquad |126-1|_5=5^{-3},
\]
so $126$ is closer to $1$ than $26$ is. Their ordinary distance tells the
opposite story. Neither ruler is malfunctioning.

Two basic calculations help prevent mistakes:
\begin{itemize}
\item $|5+25|_5=|30|_5=1/5$, the larger of $|5|_5,|25|_5$.
Terms of different sizes cannot cancel the leading digit.
\item $|5+20|_5=|25|_5=1/25$, strictly smaller than $1/5$.
Terms of equal size can cancel, just as they can over $\R$.
\end{itemize}
The product formula balances all the rulers. For $75/14$, the non-unit
factors are
\[
\frac{75}{14}\cdot2\cdot\frac13\cdot\frac1{25}\cdot7=1,
\]
coming respectively from the usual absolute value and the primes $2,3,5,7$.

% t1-pass2 examples 11a826ac
\begin{examplex}[distances from powers of $3$]
One has $|9|_3=3^{-2}=1/9$, $|6|_3=|2\cdot 3|_3=1/3$ and $|7|_3=1$. So $9$ is closer to $0$ than $6$ is, which is the opposite of the archimedean order. Also $|9+6|_3=|15|_3=1$, equal to the larger of $|9|_3$ and $|6|_3$ only after noting they differ: the larger is $1/3$, and $15=3\cdot 5$ has valuation $1$, so $|15|_3=1/3$. The sum of two numbers of different valuations has the larger valuation's absolute value.
\end{examplex}

\begin{examplex}[the product formula for $12$]
Factor $12=2^2\cdot 3$. Then $|12|_\infty=12$, $|12|_2=1/4$, $|12|_3=1/3$, and every other $|12|_p$ is $1$. The product is $12\cdot\frac14\cdot\frac13=1$. Cancelling a prime in the numerator against the $p$-adic factor is the whole identity.
\end{examplex}

\section{The \texorpdfstring{$p$}{p}-adic integers}

% t1-pass2 motive 63faee29
A $p$-adic integer is an infinite coherent string of residues modulo $p$, $p^2$, $p^3$, and so on. Ordinary integers embed as eventually constant strings in the $p$-adic sense of finite sums, and they are dense. Figure~\ref{fig:t1-minusone} shows the string for $-1$, every digit equal to $p-1$. Compactness of $\Z_p$ is the fact Chapter~9 will use whenever an infinite family of local conditions must have a simultaneous solution.

A real number is determined by its decimal approximations. A $p$-adic integer is determined by its residues.

\begin{definition}
A \emph{$p$-adic integer} is a sequence $x=(x_n)_{n\ge1}$ with $x_n\in\Z/p^n\Z$ and $x_{n+1}\equiv x_n\pmod{p^n}$ for all $n$. They form a ring $\Z_p$ under termwise addition and multiplication (the compatibility is preserved, since reduction modulo $p^n$ is a ring morphism). The integer $m$ is identified with $(m\bmod p^n)_n$; this embeds $\Z$ in $\Z_p$, since only $0$ is divisible by every $p^n$.
\end{definition}

For each $x\in\Z_p$, let $m_n\in\{0,\dots,p^n-1\}$ represent $x_n$. Compatibility says $m_{n+1}=m_n+a_np^n$ with a \emph{digit} $a_n\in\{0,\dots,p-1\}$, so $x$ is the same thing as an infinite sequence of digits $a_0,a_1,a_2,\dots$, with $m_n=a_0+a_1p+\dots+a_{n-1}p^{n-1}$. We will justify the notation $x=\sum_{i\ge0}a_ip^i$ in a moment. The \emph{valuation} $v(x)$ of $x\ne0$ is the index of its first non-zero digit, that is, the largest $v$ with $x_v=0$ (and $v=0$ if $x_1\ne0$); set $|x|_p=p^{-v(x)}$ and $|0|_p=0$. For $x\in\Z$ this is the old $v_p$, because $x_n=0$ means $p^n\mid x$.

A sequence is \emph{Cauchy} if for every $\epsilon>0$ its sufficiently
late terms have pairwise distance less than $\epsilon$; it converges to $x$
if its distance from $x$ tends to zero. A space is \emph{complete} when
every Cauchy sequence has a limit in the space. These are the same definitions
as over $\R$, with the declared metric replacing ordinary distance.

For clarity, \emph{compact} means that every cover by open sets has a finite
subcover. A set is open when each of its points has a sufficiently small
metric ball contained in it. We will prove both this covering property and
the subsequence property directly, rather than assume a topology course.

\begin{proposition}\label{prop:Zp}
\begin{enumerate}[label=(\alph*)]
\item $x\in\Z_p$ is a unit if and only if its first digit is non-zero. Every $x\ne0$ is $p^{v(x)}u$ for a unique unit $u$.
\item $\Z_p$ is an integral domain, $|xy|_p=|x|_p|y|_p$ and $|x+y|_p\le\max(|x|_p,|y|_p)$.
\item $\Z$ is dense: $|x-m_n|_p\le p^{-n}$, so $m_n\to x$, which justifies $x=\sum_ia_ip^i$.
\item $\Z_p$ is complete: every Cauchy sequence converges.
\item $\Z_p$ is compact: every sequence has a convergent subsequence.
\end{enumerate}
\end{proposition}
\begin{proof}
(a) If $xy=1$ then $x_1y_1=1$ in $\Z/p\Z$, so $x_1\ne0$. Conversely, if $x_1\ne0$, every $m_n$ is prime to $p$, hence invertible modulo $p^n$; the inverses $y_n$ are compatible (the inverse of $x_{n+1}$, reduced modulo $p^n$, is an inverse of $x_n$, and inverses are unique), and $y=(y_n)$ is the inverse of $x$. Now let $x\ne0$ with $v=v(x)$. For each $n$, the first $v$ digits vanish, so $x_{n+v}\equiv0\pmod{p^v}$ (also when $v=0$), and $m_{n+v}=p^vw_n$ with $w_n$ an integer determined modulo $p^n$; let $u_n=w_n\bmod p^n$. These are compatible, $u_1\ne0$ because $x_{v+1}\ne0$, and $(p^vu)_n=p^vw_n\bmod p^n=m_{n+v}\bmod p^n=x_n$. So $x=p^vu$ with $u$ a unit.

(b) If $u$ is a unit, $p^ku\ne0$: its residue modulo $p^{k+1}$ is $p^ku_1\ne0$; and its first non-zero digit is at index $k$, so $v(p^ku)=k$. Hence for $x=p^vu$, $y=p^wu'$ non-zero, $xy=p^{v+w}uu'\ne0$ with $uu'$ a unit, so $v(xy)=v(x)+v(y)$. In particular the decomposition in (a) is unique: $p^vu=p^wu'$ forces $v=w$ (compare valuations), then $p^v(u-u')=0$ and $u=u'$. For the inequality, if $x_n=y_n=0$ then $(x+y)_n=0$, so $v(x+y)\ge\min(v(x),v(y))$.

(c) $(x-m_n)_n=x_n-x_n=0$, so $v(x-m_n)\ge n$.

(d) Let $(x^{(k)})$ be Cauchy. For each $n$ there is $K_n$ with $|x^{(k)}-x^{(l)}|_p<p^{-n}$ for $k,l\ge K_n$, in particular $x^{(k)}_n=x^{(l)}_n$; we may take $K_1\le K_2\le\cdots$. Put $x_n=x^{(K_n)}_n$. Then $x_{n+1}=x^{(K_{n+1})}_{n+1}\equiv x^{(K_{n+1})}_n=x^{(K_n)}_n=x_n\pmod{p^n}$, so $x\in\Z_p$, and for $k\ge K_n$, $x^{(k)}_n=x_n$, i.e.\ $|x^{(k)}-x|_p\le p^{-n}$.

(e) Among the terms of a sequence, infinitely many have the same first digit; among those, infinitely many have the same second digit; and so on. Choosing the $j$-th term of the subsequence among those agreeing in the first $j$ digits (diagonal extraction) gives a subsequence that is Cauchy, hence convergent by (d).

For the covering property, suppose an open cover of $\Z_p$ has no finite
subcover. Of the $p$ first-digit classes, at least one has no finite subcover;
inside it choose a second-digit class with the same property, and continue.
The nested classes specify compatible residues, hence an element $x\in\Z_p$.
An open set of the cover containing $x$ contains a small ball around $x$.
For sufficiently large $n$, the selected class modulo $p^n$ lies in that
ball, since any two of its elements have distance at most $p^{-n}$.
One set then covers that class, contradicting its construction.
\end{proof}

A series $\sum c_n$ of $p$-adic integers converges if and only if $c_n\to0$. Indeed, the condition is necessary as for real series, and it is sufficient because, by the ultrametric inequality, $\left|\sum_{l<n\le m}c_n\right|_p\le\max_{l<n\le m}|c_n|_p$, so the partial sums are Cauchy. Thus no series whose terms tend to zero can diverge here.

\begin{example}
In $\Z_p$, $(1-p)(1+p+\dots+p^{n-1})=1-p^n\to1$, so $\displaystyle\frac1{1-p}=\sum_{i\ge0}p^i$. Multiplying by $p-1$: $-1=\sum_{i\ge0}(p-1)p^i$, the number whose digits are all $p-1$. For $p=2$: $1+2+4+8+\cdots=-1$. Your computer agrees: in two's complement arithmetic, $-1$ is stored as $\dots11111111$.
\end{example}

\subsection*{From a domain to its fraction field}

The ring $\Z_p$ has no zero divisors, but $p$ has no inverse in it.
To obtain a field, consider pairs $(a,b)$ with $a,b\in\Z_p$, $b\ne0$,
and declare
\[
(a,b)\sim(c,d)\quad\Longleftrightarrow\quad ad=bc.
\]
This is an equivalence relation. For transitivity, $ad=bc$ and $cf=de$
give $d(af-be)=0$; since $d\ne0$ in a domain, $af=be$.
The equivalence class is written $a/b$. Define
\[
\frac ab+\frac cd=\frac{ad+bc}{bd},\qquad
\frac ab\frac cd=\frac{ac}{bd}.
\]
The denominators remain nonzero. Cross-multiplication shows that replacing
either fraction by an equivalent pair leaves these answers equivalent;
for instance $ab'=a'b$ implies
$(ad+bc)b'd=(a'd+b'c)bd$.
The ring laws follow by the same finite expansions as for ordinary fractions.
The embedding $a\mapsto a/1$ is injective, and every nonzero $a/b$
has inverse $b/a$. This constructs the field $\Q_p$ without assuming
a field-of-fractions theorem.

For $a,b\ne0$, write $a=p^ru$, $b=p^sv$ with units $u,v$.
Then $a/b=p^{r-s}uv^{-1}$. Hence every nonzero $x\in\Q_p$ is uniquely
$p^ku$ with $k\in\Z$ and $u\in\Z_p^\times$.
The integer $k$ is independent of the representation: equality of fractions
gives equality of the valuations of their cross-products.
Set $v_p(x)=k$ and $|x|_p=p^{-k}$, with $v_p(0)=+\infty$.
These extend the rational valuation and absolute value.
Products add valuations, and taking a common denominator extends the
ultrametric inequality. Its equality case for unequal sizes follows too:
if $|x|>|y|$ but $|x+y|<|x|$, then
$|x|=|(x+y)-y|\le\max(|x+y|,|y|)<|x|$, impossible.

In particular
\[
\Z_p=\{x\in\Q_p:|x|_p\le1\}.
\]
Its digit expansions extend to $\Q_p$ with finitely many negative powers:
multiply by a sufficiently large power of $p$ to enter $\Z_p$, expand,
then shift back. For example $1/p$ is in $\Q_p$ and has valuation $-1$;
its first residue modulo $p$ as a $p$-adic integer does not exist.

\subsection*{Density and completeness of the field}

For $x=p^{-k}z$ with $z\in\Z_p$, $k\ge0$, let $m_n$ be the first
$n$ digits of $z$. Then $p^{-k}m_n\in\Q$ and
$|x-p^{-k}m_n|_p\le p^{k-n}\to0$. Thus $\Q$ is dense in $\Q_p$.
This also explains precisely why $\Z$ is dense in $\Z_p$.

Let $(x_j)$ be Cauchy in $\Q_p$. It is bounded:
after some index $J$, $|x_j-x_J|_p\le1$, so
$|x_j|_p\le\max(1,|x_J|_p)$ there; the finitely many earlier sizes
can be included in a common bound. Choose $k\ge0$ with
$|p^kx_j|_p\le1$ for all $j$. This scaled Cauchy sequence lies in
the complete ring $\Z_p$ and has a limit $y$ there.
Therefore $x_j\to p^{-k}y$ in $\Q_p$.
We have proved completeness, not just adjoined a name to formal digits.

In this complete field, the same ultrametric tail estimate proves
$\sum c_n$ converges if and only if $c_n\to0$.
The ordinary harmonic series $\sum_{n\ge1}1/n$ still diverges:
$|1/n|_p\ge1$ for every positive integer $n$, so its terms do not vanish.

\subsection*{Arithmetic with a declared precision}
Writing $x=\cdots+a_2p^2+a_1p+a_0$ is a compact way of storing compatible
remainders. At precision $p^N$, we simply calculate in $\Z/p^N\Z$.
For instance, in $\Z_5$,
\[
\frac13\equiv42=2+3\cdot5+1\cdot25\pmod{125},
\]
because $3\cdot42=126\equiv1\pmod{125}$. This says
$|1/3-42|_5\le5^{-3}$, not that $42$ approximates $1/3$ in the usual sense.

Addition and multiplication by integral elements preserve an error
$O(p^N)$, meaning an element of $p^N\Z_p$. Division by a unit preserves it too.
Division by $p$ loses one known digit:
\[
x\equiv y\pmod{p^N}\quad\Longrightarrow\quad
x/p\equiv y/p\pmod{p^{N-1}}.
\]
This is the bookkeeping behind Newton computations when the derivative is
not a unit. Integer congruences are exact; finite precision does not mean
floating-point rounding.

\subsection*{Why compactness is useful}
Proposition~\ref{prop:Zp}(e) proves both compactness and the convergent-subsequence property directly. Here is a consequence of the latter.
Suppose $f\in\Z_p[X]$ has a root modulo $p^n$ for every $n$, with no promise
that the roots you picked are compatible. Choose a representative $b_n\in\Z$
of each root. A convergent subsequence $b_{n_j}\to b\in\Z_p$ exists. Since
$|f(b_{n_j})|_p\le p^{-n_j}\to0$ and polynomials are continuous, $f(b)=0$.
Thus roots at every finite level yield a genuine $p$-adic root. A root at
one level alone makes no such promise.

\begin{proposition}\label{prop:periodic}
A $p$-adic number is rational if and only if its digit sequence is eventually periodic.
\end{proposition}
\begin{proof}
If the digits are eventually periodic with period $\ell$, then, after subtracting a rational number (finitely many terms) and dividing by a power of $p$, we are left with $x=c\,(1+p^\ell+p^{2\ell}+\cdots)=\frac c{1-p^\ell}$ with $c\in\Z$: rational. Conversely, let $x=\frac ab$ be rational; multiplying by a power of $p$ (which shifts the digits), we may assume $p\nmid b$, $b\ge1$. The first digit $a_0$ is the residue of $ab^{-1}$ modulo $p$, and $\frac{x-a_0}p=\frac{a'}b$ with $a'=\frac{a-a_0b}p\in\Z$: the same denominator. Iterating produces the digits from a sequence of numerators $a,a',a'',\dots$ with
\[
|a'|\le\frac{|a|+(p-1)b}p.
\]
If $|a|>b$, this gives $|a'|<\frac{|a|+(p-1)|a|}p=|a|$; if $|a|\le b$, it gives $|a'|\le b$. So the numerators decrease until they enter $[-b,b]$ and then stay there. A finite set of states, each determining the next one, must eventually cycle, so the digits are eventually periodic.
\end{proof}

% t1-pass2 examples 63faee29
\begin{examplex}[the $5$-adic expansion of $-1$]
The geometric series $\sum_{k\ge 0}4\cdot 5^k$ has partial sums $4$, $24$, $124$, $624$. Each is $-1$ modulo the next power of $5$: $4\equiv -1\pmod 5$, $24\equiv -1\pmod{25}$, $124\equiv -1\pmod{125}$. So the digit string $\ldots 444_5$ is $-1$ in $\Z_5$. The same string with digit $p-1$ is $-1$ in every $\Z_p$.
\end{examplex}

\begin{figure}[ht]
\centering
\fig{ch13-minusone}
\caption{The first three digits of $-1$ in $\Z_5$, and the partial sums, each of which is already $-1$ at that precision. Every further digit is again $4$.}
\label{fig:t1-minusone}
\end{figure}

\begin{examplex}[a unit and a non-unit]
The $5$-adic integer with first digit $2$ is a unit, because $2$ is invertible modulo $5$ and Hensel's lemma, or simply the explicit inverse of each partial sum, lifts the inverse. The integer with first digit $0$ is divisible by $5$ and is not a unit. In $\Z_p$ one sees units at the first digit.
\end{examplex}

\section{Hensel's lemma, Newton's way}

% t1-pass2 motive ac26f553
Newton's method is an algorithm for lifting a simple root modulo $p$ to a root in $\Z_p$. The derivative must be a unit, or at least large compared with the value, or precision is lost. This is the tool that constructs $\sqrt{-1}$ in $\Z_5$ and decides the squares in $\Q_p$. Chapter~9 will apply the same Newton step in every complete local field.

Volume~2 lifted a root modulo $p$ with $f'(a)\not\equiv0\pmod p$ to every $p^k$. In $\Z_p$ the lifts converge to an actual root, and Newton's method does it with quadratic speed, even in some cases where $f'(a)\equiv0\pmod p$.

\begin{think}
$x^2+7$ has $x=1$ as a root modulo $8$, but $f'(1)=2$ is not a unit in $\Z_2$, so Volume~2's lemma does not apply. Does $-7$ have a square root in $\Z_2$? Try to find one modulo $16$, $32$, $64$.
\end{think}

\begin{theorem}[Hensel's lemma, Newton's form]\label{thm:henselnewton}
Let $f\in\Z_p[X]$ and $a\in\Z_p$ with $|f(a)|_p<|f'(a)|_p^2$. Then there is a unique $\alpha\in\Z_p$ with $f(\alpha)=0$ and $|\alpha-a|_p<|f'(a)|_p$. Moreover $|\alpha-a|_p\le\left|\frac{f(a)}{f'(a)}\right|_p$.
\end{theorem}
\begin{proof}
Write $|\cdot|$ for $|\cdot|_p$. Taylor's formula for polynomials gives $f(x+h)=f(x)+f'(x)h+h^2g(x,h)$ and $f'(x+h)=f'(x)+h\,k(x,h)$, with $g,k$ polynomials with coefficients in $\Z_p$ (the coefficients of $\frac{f^{(j)}}{j!}$ are binomial multiples of those of $f$). So for $x,h\in\Z_p$: $|g(x,h)|,|k(x,h)|\le1$.

Let $\delta=|f'(a)|>0$ and $t=\frac{|f(a)|}{\delta^2}<1$. Newton's method sets $a_0=a$ and $a_{j+1}=a_j+h_j$ with $h_j=-\frac{f(a_j)}{f'(a_j)}$. We show by induction that $a_j\in\Z_p$, $|f'(a_j)|=\delta$ and $|f(a_j)|\le\delta^2t^{2^j}$. For $j=0$ this is the hypothesis. Given it for $j$: $|h_j|\le\delta t^{2^j}<\delta\le1$, so $a_{j+1}\in\Z_p$; $f'(a_{j+1})=f'(a_j)+h_jk$ with $|h_jk|<\delta=|f'(a_j)|$, so $|f'(a_{j+1})|=\delta$ (Proposition~\ref{prop:ultra}(b)); and $f(a_{j+1})=f(a_j)+f'(a_j)h_j+h_j^2g=h_j^2g$, so $|f(a_{j+1})|\le|h_j|^2\le\delta^2t^{2^{j+1}}$.

Hence $|a_{j+1}-a_j|=|h_j|\le\delta t^{2^j}\to0$, so $(a_j)$ is Cauchy and converges to some $\alpha\in\Z_p$ (Proposition~\ref{prop:Zp}(d)), and $f(\alpha)=\lim f(a_j)=0$ since polynomials are continuous. Also $|\alpha-a|\le\max_j|h_j|\le\delta t=\left|\frac{f(a)}{f'(a)}\right|<\delta$.

Uniqueness: let $f(\beta)=0$ with $|\beta-a|<\delta$, and $h=\beta-\alpha$, so $|h|<\delta$. As above $|f'(\alpha)|=\delta$. Then $0=f(\beta)=f(\alpha)+f'(\alpha)h+h^2g(\alpha,h)=h\big(f'(\alpha)+hg(\alpha,h)\big)$, and the bracket is non-zero because $|hg|<\delta=|f'(\alpha)|$. So $h=0$.
\end{proof}

% t1-pass2 why
\paragraph{Why this works.}
The correction $-f(a)/f'(a)$ is smaller than the previous error precisely when $|f(a)|<|f'(a)|^2$. Ultrametric geometry keeps the derivative from changing at the new point, so the same inequality persists and the errors form a geometric sequence with ratio $t<1$. Completeness of $\Z_p$ supplies the limit, and a second root in the same ball would force the derivative to vanish.

The proof shows that the number of correct digits roughly doubles at each step, exactly as for Newton's method over $\R$.

\subsection*{A square root appears one digit at a time}
Take $f(X)=X^2+1$ in $\Z_5[X]$, starting at $a=2$. Here $f(2)=5$ and
$f'(2)=4$ is a unit. Instead of rational Newton iterates, we can lift one
base-$5$ digit at a time. If $b$ is a root modulo $5^k$, write
$b'=b+t5^k$ with $t\in\{0,1,2,3,4\}$. The condition at the next level is
\[
\frac{b^2+1}{5^k}+2bt\equiv0\pmod5.
\]
Because $2b$ is a unit, this determines exactly one digit $t$. Starting with
$b=2$, the successive representatives are $2,7,57,182$ modulo
$5,25,125,625$. The corresponding digits are $2,1,2,1$.

\begin{figure}[ht]
\centering\fig{ch12-lifting}
\caption{The root of $X^2+1$ congruent to $2$ modulo $5$. At each level,
only one of the five candidate lifts survives; the solid path records
$2,7,57,182$. Dashed branches show the four rejected choices at each step.}
\label{fig:hensel-lifting}
\end{figure}

For $p=2$ a caution is needed: $X^2+7$ has four roots modulo $256$, namely
$53,75,181,203$, but only two roots in the field $\Q_2$. They have residues
$75,181$ modulo $256$. The other finite roots are not the residues of actual
$2$-adic roots: they already fail to lift to roots modulo $512$.
A polynomial over a field has at most its degree many roots; a polynomial
over a ring with zero divisors, such as $\Z/256\Z$, can have more.

\begin{think}
Why does checking $a^2+7\equiv0\pmod{256}$ not certify eight correct digits
of an actual square root? The derivative $2a$ is divisible by $2$, so solving
for the correction loses precision. Here Hensel's bound only certifies
$|\alpha-a|_2\le|f(a)/f'(a)|_2\le2^{-7}$ unless you know an extra digit
of $f(a)$. This is a practical reason to read the error bound as well as the
existence statement.
\end{think}

\begin{corollary}[squares in $\Q_p$]\label{cor:squares}
A non-zero $x=p^vu\in\Q_p$ ($u$ a unit) is a square if and only if $v$ is even and
\begin{itemize}
\item for $p$ odd: $u\bmod p$ is a square modulo $p$;
\item for $p=2$: $u\equiv1\pmod8$.
\end{itemize}
In particular $\sqrt{-7}\in\Z_2$, and $\sqrt{-1}\in\Z_p$ if and only if $p\equiv1\pmod4$.
\end{corollary}
\begin{proof}
If $x=y^2$ with $y=p^wu'$, then $v=2w$ and $u=u'^2$, so $u\bmod p$ is a square, and for $p=2$, $u'$ is odd and odd squares are $\equiv1\pmod8$ ($(2k+1)^2=4k(k+1)+1$ with $k(k+1)$ even). Conversely, it suffices to show $u$ is a square. For $p$ odd, choose $b\in\Z$ with $b^2\equiv u\pmod p$ and $f=X^2-u$: $|f(b)|\le p^{-1}<1=|2b|^2=|f'(b)|^2$. For $p=2$ and $u\equiv1\pmod 8$, take $b=1$: $|f(1)|=|1-u|\le2^{-3}<2^{-2}=|f'(1)|^2$. In both cases Theorem~\ref{thm:henselnewton} gives a root. The last claims: $-7\equiv1\pmod8$, and $-1$ is a square modulo $p$ if and only if $p\equiv1\pmod4$ (Volume~2); for $p=2$, $-1\not\equiv1\pmod8$.
\end{proof}

% t1-pass2 why
\paragraph{Why this works.}
A square has even valuation, and the unit part must be a square in $\Z_p$. For odd $p$ that is the Legendre symbol, lifted by Hensel because the derivative $2b$ is a unit. For $p=2$ the first modulus at which odd squares are distinguished is $8$, and $u\equiv 1\pmod 8$ is exactly the condition that makes $|f(1)|<|f'(1)|^2$ for $f=X^2-u$.

\begin{corollary}[roots of unity]\label{cor:teich}
$\Z_p$ contains exactly $p-1$ roots of the polynomial $X^{p-1}-1$, one congruent to each $a\in\{1,\dots,p-1\}$ modulo $p$.
\end{corollary}
\begin{proof}
For $f=X^{p-1}-1$ and $a\in\{1,\dots,p-1\}$: $f(a)\equiv0\pmod p$ by Fermat, and $f'(a)=(p-1)a^{p-2}$ is a unit, so Theorem~\ref{thm:henselnewton} gives a root $\omega(a)\equiv a\pmod p$, the unique one in $B(a,p^{-1})$. Roots in different classes are different, and a polynomial of degree $p-1$ over the integral domain $\Z_p$ has at most $p-1$ roots (the argument of Volume~1, the polynomial root-bound proposition, works in any integral domain).
\end{proof}

So $\Q_5$ contains $i$ and $\Q_7$ contains $e^{2i\pi/3}$, in the algebraic sense: roots of $X^2+1$ and $X^2+X+1$. The numbers $\omega(a)$ are the \emph{Teichm\"uller representatives}; Exercise~\ref{ex:teich} computes them as limits.

% t1-pass2 examples ac26f553
\begin{examplex}[lifting a square root of $-1$ modulo $5$]
Modulo $5$, $2^2=4\equiv -1$ and the derivative of $X^2+1$ at $2$ is $4\not\equiv 0$. Newton's step from $a=2$: $f(2)=5$, so the correction is $-5\cdot(4)^{-1}$. Modulo $25$ one looks for $2+5t$ with $(2+5t)^2\equiv -1\pmod{25}$, that is $4+20t\equiv -1\pmod{25}$, so $20t\equiv -5\pmod{25}$. Dividing by $5$ gives $4t\equiv -1\pmod 5$, so $t\equiv 1$ and the lift is $7$. Indeed $7^2=49\equiv -1\pmod{25}$.
\end{examplex}

\begin{examplex}[a root that does not lift]
The polynomial $X^2-5$ has a root modulo $5$, namely $0$, but $f'(0)=0$. Modulo $25$ the congruence $x^2\equiv 5\pmod{25}$ has no solution: squares of multiples of $5$ are $0$ modulo $25$, and other squares are not $5$. The derivative hypothesis is the difference between this failure and the success for $X^2+1$.
\end{examplex}

\section{Ostrowski's theorem}

% t1-pass2 motive 89f845ac
One might invent other ways to measure rational numbers. Ostrowski's theorem says there are none, up to equivalence: every absolute value on $\Q$ is a power of the usual one or of some $|\cdot|_p$. The list of places in Chapter~9 is therefore a theorem, not a convention.

An \emph{absolute value} on $\Q$ is a map $\|\cdot\|:\Q\to[0,+\infty)$ with $\|x\|=0\iff x=0$, $\|xy\|=\|x\|\,\|y\|$ and $\|x+y\|\le\|x\|+\|y\|$. Examples: the \emph{trivial} one ($\|x\|=1$ for $x\ne0$), $|\cdot|_\infty$, the $|\cdot|_p$, and their powers $|\cdot|_\infty^s$ ($0<s\le1$) and $|\cdot|_p^s$ ($s>0$), which define the same notion of convergence.

\begin{think}
Is there any other natural way to say that two rational numbers are close? Before reading on, try to invent one, say from the digits in base $10$, and check whether it is multiplicative.
\end{think}

\begin{theorem}[Ostrowski, 1916]\label{thm:ostrowski}
Every non-trivial absolute value on $\Q$ is $|\cdot|_\infty^s$ for some $0<s\le1$, or $|\cdot|_p^s$ for a prime $p$ and some $s>0$.
\end{theorem}
\begin{proof}
First, $\|1\|=\|1\|^2$ gives $\|1\|=1$, and $\|-1\|^2=1$ gives $\|-1\|=1$, so $\|-x\|=\|x\|$; by multiplicativity it suffices to know $\|n\|$ for integers $n\ge2$.

\emph{Case 1: $\|n_0\|>1$ for some integer $n_0\ge2$.} Let $m,n\ge2$. Write $m^k$ in base $n$: $m^k=\sum_{i=0}^Lc_in^i$ with $0\le c_i\le n-1$ and $L\le k\frac{\log m}{\log n}$. Then
\[
\begin{aligned}&\|m\|^k\le\sum_{i=0}^L\|c_i\|\,\|n\|^i\eqbr\le(L+1)(n-1)\max(1,\|n\|)^L,\end{aligned}
\]
since $\|c\|\le c$ for integers $c\ge0$ (triangle inequality). Taking $k$-th roots and letting $k\to\infty$, $(L+1)^{1/k}\to1$ and
\[
\|m\|\le\max(1,\|n\|)^{\log m/\log n}.\qquad(\ast)
\]
If some $\|n\|\le1$, $(\ast)$ gives $\|m\|\le1$ for every $m$, contradicting $\|n_0\|>1$. So $\|n\|>1$ for all $n\ge2$, and $(\ast)$ reads $\|m\|^{1/\log m}\le\|n\|^{1/\log n}$. Exchanging $m,n$, these are equal: $\|n\|=n^s$ for a constant $s>0$, that is, $\|x\|=|x|_\infty^s$. Finally $2^s=\|2\|\le\|1\|+\|1\|=2$ forces $s\le1$.

\emph{Case 2: $\|n\|\le1$ for all integers $n$.} Then $\|\cdot\|$ is ultrametric: for $x,y\in\Q$ and $k\ge1$,
\[
\begin{aligned}\|x+y\|^k&=\Big\|\sum_{i=0}^k\binom kix^iy^{k-i}\Big\|\eqbr\le(k+1)\max(\|x\|,\|y\|)^k,\end{aligned}
\]
and $k$-th roots give $\|x+y\|\le\max(\|x\|,\|y\|)$. Since the absolute value is non-trivial, some integer $n\ge2$ has $\|n\|<1$ (otherwise all non-zero integers, hence all non-zero rationals, have absolute value $1$). The set $I=\{n\in\Z:\|n\|<1\}$ is then an ideal of $\Z$ (stable under sums by the ultrametric inequality, and under multiplication by integers since $\|m\|\le1$), different from $\Z$ ($\|1\|=1$) and from $\{0\}$. It is prime: if $\|ab\|<1$ then $\|a\|<1$ or $\|b\|<1$, as both are $\le1$. So $I=p\Z$ for a prime $p$ (Volume~2). Now write $n=p^vm$ with $p\nmid m$: $\|m\|=1$, so $\|n\|=\|p\|^v=(p^{-v})^s$ with $s=-\frac{\log\|p\|}{\log p}>0$. Hence $\|x\|=|x|_p^s$.
\end{proof}

% t1-pass2 why
\paragraph{Why this works.}
If some integer has absolute value greater than $1$, the axioms force the absolute value to grow like a power of the usual size. If every integer has absolute value at most $1$, then some prime is strictly smaller than $1$, the ultrametric inequality follows by expanding binomial coefficients, and every rational is a unit times a power of that prime. Those are the only two shapes.

Consequently, up to equivalence of non-trivial absolute values, the completions of $\Q$ are $\R$ and the $\Q_p$. (For the trivial absolute value, $\Q$ is already complete with the discrete metric.) The equivalence classes of non-trivial absolute values are the \emph{places} of $\Q$; their associated completions are these fields, and the product formula of Proposition~\ref{prop:ultra} ties them together.

% t1-pass2 examples 89f845ac
\begin{examplex}[a false absolute value]
Declare two rationals close when their difference has a large power of $10$ in the denominator of its decimal expansion, without requiring multiplicativity. The resulting notion of closeness need not satisfy $|xy|=|x||y|$. Ostrowski's first move is to derive, from the axioms, that the value of every integer is governed by a single prime or by size. A definition that is not multiplicative never enters the proof.
\end{examplex}

\begin{examplex}[powers do not give new places]
The function $|x|_5^2$ is an absolute value, since the exponent $2$ preserves the ultrametric inequality. It determines the same convergent sequences as $|\cdot|_5$. The theorem classifies absolute values up to this reparametrisation: the exponent $s$ in $|x|_p^s$ is not a new prime.
\end{examplex}

\section{Local and global}

% t1-pass2 motive 1e195763
A conic can have points over $\R$ and over every $\Q_p$ and still, for cubics, fail over $\Q$. For quadrics the local conditions are sufficient: that is the Hasse--Minkowski theorem, proved in Chapter~10. This section gives the elementary half, an obstruction that is visible in $\Q_3$ and therefore visible in $\Q$.

\begin{think}
If an equation with integer coefficients has a rational solution, it has solutions in $\R$ and in every $\Q_p$. What about the converse? Look at $x^2+y^2=3z^2$ modulo $9$ before reading on.
\end{think}

A rational solution is a \emph{global} solution, a solution in $\R$ or $\Q_p$ a \emph{local} one. Local tests are often much more manageable: signs and continuity help over $\R$, while congruences and Hensel's lemma help over $\Q_p$. A successful Hensel test proves existence, but a failed test alone does not prove non-existence, especially at a singular root. The \emph{Hasse principle} asks whether local solvability everywhere implies global solvability.

\begin{proposition}\label{prop:noconic}
The equation $x^2+y^2=3z^2$ has no solution in $\Q_3$ other than $(0,0,0)$. Hence it has no non-zero rational solution, although it has real ones.
\end{proposition}
\begin{proof}
Suppose $(x,y,z)\ne0$ in $\Q_3^3$. Multiplying by a power of $3$, we may assume $x,y,z\in\Z_3$, not all divisible by $3$. Squares modulo $3$ are $0,1$, so $x^2+y^2\equiv0\pmod 3$ forces $3\mid x$ and $3\mid y$. Then $9\mid3z^2$, so $3\mid z$: a contradiction.
\end{proof}

% t1-pass2 why
\paragraph{Why this works.}
Squares modulo $3$ are $0$ and $1$, so a sum of two squares is $0$ modulo $3$ only when both are $0$. The equation forces that vanishing, hence $3$ divides every variable, and the solution was not primitive. An infinite descent on the power of $3$ remains. The real place never sees this obstruction.

\subsection*{What ``primitive'' contributes to a local obstruction}
A non-zero vector in $\Q_p^3$ can be multiplied by a power of $p$ until all
coordinates are integral and at least one is a unit: take the negative of the
smallest coordinate valuation. Such a vector is called \emph{primitive}.
For a homogeneous equation, scaling preserves the equation. Thus to rule out
non-zero $p$-adic solutions it suffices to rule out primitive integral ones.

For $x^2+y^2=3z^2$, the contradiction is a two-stage test: modulo $3$ we get
$3\mid x,y$; modulo $9$ we then get $3\mid z$. It does \emph{not} say that the
only solution modulo $9$ is $(0,0,0)$: $(3,0,0)$ is a non-zero solution there.
It says every such solution is non-primitive. Distinguishing those statements
is the difference between a useful obstruction and a false one.

A positive example is $x^2+y^2=5z^2$, with the rational point $(1,2,1)$.
It automatically gives a point in every completion. The hard direction of a
local--global theorem goes the other way: independently obtained local points
need not be the shadows of one rational point.

\begin{remark}
For quadratic forms the converse holds: by the \emph{Hasse--Minkowski theorem} (1921), a quadratic form with rational coefficients represents $0$ non-trivially over $\Q$ if and only if it does so over $\R$ and over every $\Q_p$. It is proved in Chapter~10, and it decides every conic by finitely many congruence checks. For cubics it fails: Selmer (1951) showed that $3x^3+4y^3+5z^3=0$ has non-zero solutions in $\R$ and in every $\Q_p$, but none in $\Q$. Measuring this failure (the Tate--Shafarevich group) is one of the deepest themes of modern number theory.
\end{remark}

% t1-pass2 examples 1e195763
\begin{examplex}[$x^2+y^2=3$ over the reals and over $\Q_3$]
Real points exist. In $\Q_3$, a primitive solution of $X^2+Y^2=3Z^2$ would have $X$ and $Y$ not both divisible by $3$, but squares modulo $3$ are $0$ or $1$, so the left side is $0,1$ or $2$ modulo $3$ and equals $0$ only if both squares vanish. Thus $3$ divides $X$ and $Y$, then $3$ divides $Z$. There is no primitive local solution, hence no rational solution.
\end{examplex}

\begin{examplex}[a square that exists locally everywhere and globally]
The integer $2$ is a square in $\R$ and, by the criterion of this chapter, one can ask in which $\Q_p$ it is a square. It is a square in $\Q_7$ because $\bigl(\frac{2}{7}\bigr)=1$ ($7\equiv \pm 1\pmod 8$? $7\equiv -1\pmod 8$, yes). Globally $2$ is not a square in $\Q$. There is no contradiction: being a square in $\Q$ demands a rational root of $X^2-2$, which is a stronger equation than the local solubility of a conic. The Hasse principle is a theorem about quadratic forms, not about every polynomial.
\end{examplex}

\begin{beauty}
Where this leads. Chapter~9 takes $\Q_p$ as the local field at $p$, Chapter~10 proves the Hasse--Minkowski theorem, and Chapter~12 builds Tate's thesis on the same absolute values. The square classes computed in the exercises are the local input of class field theory. Ostrowski's list is the reason there is one Euler factor for each prime and one for the real place, and no others.
\end{beauty}

\section{Exercises}

\begin{exercise}[$\star$]\label{ex:padicdigits}
Find the $2$-adic digits of $\frac13$ and of $-\frac13$, and the $5$-adic digits of $\frac14$. Check Proposition~\ref{prop:periodic} on these examples.
\end{exercise}

\begin{exercise}[$\star$]\label{ex:nfact}
Show that $\sum_{n\ge0}n!$ converges in every $\Z_p$, and that $\sum_{n\ge0}n\cdot n!=-1$ in every $\Z_p$.
\end{exercise}

\begin{exercise}[$\star$]\label{ex:sqrtm7}
Compute the last $8$ binary digits of the two square roots of $-7$ in $\Z_2$, with Newton's method started at $a=1$ and at $a=3$.
\end{exercise}

For the next exercise, a square class identifies nonzero $x,y\in\Q_p$
when $x/y$ is a square. This is an equivalence relation, compatible with
multiplication, since products and inverses of nonzero squares are squares.
The notation $\Q_p^\times/(\Q_p^\times)^2$ means precisely these classes.

\begin{exercise}[$\star\star$]\label{ex:sqclasses}
Which of $-1,2,3,5,-7$ are squares in $\Q_2$, in $\Q_3$, in $\Q_5$? Show that $\Q_p^\times/(\Q_p^\times)^2$ has $4$ elements for $p$ odd and $8$ for $p=2$.
\end{exercise}

\begin{exercise}[$\star\star$]\label{ex:padicexp}
Using Legendre's formula for $v_p(n!)$ (Volume~2), show that $v_p(n!)=\frac{n-s_p(n)}{p-1}$, where $s_p(n)$ is the sum of the base-$p$ digits of $n$. Deduce that $\exp(x)=\sum_{n\ge0}\frac{x^n}{n!}$ converges in $\Q_p$ if and only if $|x|_p<p^{-1/(p-1)}$. Where does it converge in $\Q_2$?
\end{exercise}

\begin{exercise}[$\star\star$]\label{ex:hasse}
Show that $x^2-2y^2=3z^2$ has no non-zero rational solution, and that $x^2+y^2=7z^2$ has none either. Which prime is responsible in each case?
\end{exercise}

\begin{exercise}[$\star\star\star$]\label{ex:teich}
Let $a\in\Z$ with $p\nmid a$. Show that if $x\equiv y\pmod{p^k}$ with $k\ge1$, then $x^p\equiv y^p\pmod{p^{k+1}}$. Deduce that $(a^{p^n})_{n\ge0}$ converges in $\Z_p$, and that its limit is the Teichm\"uller representative $\omega(a)$ of Corollary~\ref{cor:teich}. Compute $\omega(2)$ in $\Z_5$ modulo $125$.
\end{exercise}

\begin{exercise}[$\star\star$]\label{ex:t1-random-valuation}
\emph{A finite model of random divisibility.} Fix a prime $p$ and $N\ge1$.
Choose $U_N$ uniformly from $\{0,\dots,p^N-1\}$, and put
$V_N=\min(v_p(U_N),N)$, with $v_p(0)=+\infty$.
Find $\operatorname{Pr}(V_N\ge k)$ for $0\le k\le N$, the full distribution
and expectation of $V_N$. What distribution do the individual masses
approach as $N$ grows? Explain why this is not a uniform random positive integer.
\end{exercise}

\begin{exercise}[$\star$]\label{ex:t1-minusone}
Show that $\sum_{k\ge 0}4\cdot 5^k=-1$ in $\Z_5$, by checking that the partial sum of the first $n$ terms is congruent to $-1$ modulo $5^n$ for $n=1,2,3$.
\end{exercise}

\begin{beauty}
The $p$-adic numbers turn congruences into analysis: Hensel's lemma is Newton's method, and the compactness of $\Z_p$ is the reason why a congruence solvable modulo every $p^k$ has a $p$-adic solution. Ostrowski's theorem says that $\R$ is just one of infinitely many equally legitimate completions of $\Q$, and the modern point of view treats them all at once, in the ring of \emph{ad\`eles} (Chapter~12). The $p$-adics are also everywhere in computing: two's complement integers are $2$-adic numbers truncated to $64$ digits, and exact linear algebra and polynomial factorisation work modulo $p^k$ and lift. And deep results live here: the Skolem--Mahler--Lech theorem on zeros of linear recurrences and Kummer's congruences for Bernoulli numbers are both proved by $p$-adic analysis.
\end{beauty}

\section{Hints}

\begin{hint}{ex:padicdigits}
For $\frac13$ in $\Z_2$, solve $3x\equiv1$ modulo $2,4,8,16$, or use $\frac13=\frac{-1}{1-4}$.
\end{hint}
\begin{hint}{ex:nfact}
A $p$-adic series converges if and only if its terms tend to $0$. For the second sum, telescope: $n\cdot n!=(n+1)!-n!$.
\end{hint}
\begin{hint}{ex:sqrtm7}
Newton's step for $X^2+7$ is $a\mapsto\frac{a^2-7}{2a}$; compute modulo $2^{10}$ to keep track of the precision lost in the division by $2$.
\end{hint}
\begin{hint}{ex:sqclasses}
Use Corollary~\ref{cor:squares}. A square class is determined by the parity of $v$ and the class of $u$ modulo squares.
\end{hint}
\begin{hint}{ex:padicexp}
Write $n=\sum c_ip^i$ and compute $\sum_k\lfloor n/p^k\rfloor$ digit by digit. Then $v_p\left(\frac{x^n}{n!}\right)=n\,v_p(x)-\frac{n-s_p(n)}{p-1}$.
\end{hint}
\begin{hint}{ex:hasse}
Find a prime $p$ and a power $p^k$ for which every solution modulo $p^k$ has all coordinates divisible by $p$, as in Proposition~\ref{prop:noconic}. This rules out primitive solutions, rather than every non-zero residue vector.
\end{hint}
\begin{hint}{ex:teich}
For the lemma, write $x=y+tp^k$ and expand with the binomial theorem. Then $a^{p^{n+1}}$ and $a^{p^n}$ are congruent modulo $p^{n+1}$.
\end{hint}

\begin{hint}{ex:t1-random-valuation}
Count the $p^{N-k}$ residues divisible by $p^k$.
Use $\mathbb E V_N=\sum_{k=1}^N\operatorname{Pr}(V_N\ge k)$,
checking it first for each integer value of $V_N$.
\end{hint}

\begin{hint}{ex:t1-minusone}
The partial sum is $4(5^n-1)/(5-1)=5^n-1$.
\end{hint}

\section{Solutions}

\begin{solution}{ex:padicdigits}
\idea Two methods. Either write the fraction as $\frac c{1-p^\ell}$, whose expansion repeats the digits of $c$ every $\ell$ places; or compute digits one at a time: the first digit $a_0$ is $x\bmod p$, then replace $x$ by $\frac{x-a_0}p$ and repeat, as in the proof of Proposition~\ref{prop:periodic}.
\writeup $-\frac13=\frac1{1-4}=1+4+16+64+\cdots$, so the $2$-adic digits of $-\frac13$ (from $2^0$ upward) are $1,0,1,0,1,0,\dots$. Then $\frac13=1+2\cdot\left(-\frac13\right)=1+2+8+32+\cdots$, with digits $1,1,0,1,0,1,\dots$. Check: $1+2+8+32=43$ and $3\cdot43=129\equiv1\pmod{64}$. For $\frac14$ in $\Z_5$, digit by digit: $\frac14\equiv4\pmod5$ (as $4\cdot4=16\equiv1$), and $\frac{1/4-4}5=-\frac34$. Next $-\frac34\equiv3\pmod5$ (as $4\cdot3=12\equiv-3$), and $\frac{-3/4-3}5=-\frac34$ again. So the digits are $4,3,3,3,\dots$, i.e.\ $\frac14=4+3\cdot5+3\cdot25+\cdots$. Check: $4+15+75=94$ and $4\cdot94=376=3\cdot125+1$. All three expansions are eventually periodic, as Proposition~\ref{prop:periodic} predicts.
\end{solution}

\begin{solution}{ex:nfact}
\idea In $\Z_p$, $n!\to0$ because $n!$ is divisible by more and more factors $p$. The second sum telescopes, and its partial sums are $(N+1)!-1$.
\writeup By Legendre's formula, $v_p(n!)\ge\lfloor n/p\rfloor\to\infty$, so $|n!|_p\to0$ and $\sum n!$ converges. Next, $n\cdot n!=(n+1)!-n!$, so $\sum_{n=0}^Nn\cdot n!=(N+1)!-0!=(N+1)!-1$, which tends to $-1$ in $\Z_p$ because $(N+1)!\to0$.
\end{solution}

\begin{solution}{ex:sqrtm7}
\idea Run Newton's iteration $a\mapsto a-\frac{a^2+7}{2a}=\frac{a^2-7}{2a}$ in exact rational arithmetic (the denominators stay odd), and reduce modulo $2^8$ only at the end. The proof of Theorem~\ref{thm:henselnewton} tells us how many steps are needed.
\writeup For $f=X^2+7$ and $a=1$: $|f'(1)|_2=\frac12$ and $|f(1)|_2=|8|_2=\frac18<\frac14$, so the theorem applies with $t=\frac12$, and $|f(a_j)|_2\le\frac14\cdot2^{-2^j}$. From $a_0=1$: $a_1=-3$, $a_2=-\frac13$, $a_3=\frac{31}3$, $a_4=\frac{449}{93}$. Indeed $a_4^2+7=\frac{262144}{8649}=\frac{2^{18}}{93^2}$, so $a_4$ is a square root of $-7$ modulo $2^{18}$, hence modulo $2^8$ the root is $449\cdot93^{-1}\equiv181\pmod{256}$; check: $181^2+7=32768=2^{15}$. In binary, $181=10110101$. The other root is $-181\equiv75=01001011$ in binary, and Newton's method started at $a_0=3$ reaches it: $3,\frac13,-\frac{31}3,-\frac{449}{93},\dots$ Note how the two roots are negatives of each other: their binary expansions are two's complements.
\end{solution}

\begin{solution}{ex:sqclasses}
\idea Corollary~\ref{cor:squares} reduces everything to a congruence on the unit part: modulo $p$ for odd $p$, modulo $8$ for $p=2$. Counting square classes is then counting the possible (parity of $v$, class of $u$) pairs.
\writeup In $\Q_2$: $-1\equiv7$, $3$, $5\pmod8$ are not $1$, and $2$ has odd valuation, so none of $-1,2,3,5$ is a square, while $-7\equiv1\pmod8$ is. In $\Q_3$: the squares modulo $3$ are $\{1\}$; $-1\equiv2$ and $5\equiv2$ are not squares, $-7\equiv2$ neither, $2$ is not, and $3$ has odd valuation; none is a square. In $\Q_5$: squares modulo $5$ are $\{1,4\}$; $-1\equiv4$ is a square, $2$ and $3$ are not, $5$ has odd valuation, and $-7\equiv3$ is not. For the count with $p$ odd: $x=p^vu$ is determined modulo squares by $v\bmod2$ and by whether $u\bmod p$ is a square, and Corollary~\ref{cor:squares} shows these two bits are independent and both occur: $4$ classes, represented by $1,\varepsilon,p,\varepsilon p$ with $\varepsilon$ a non-square unit. For $p=2$: $v\bmod2$, and the class of $u$ modulo $8$ among $1,3,5,7$, since $u$ is a square times $u'$ if and only if $u\equiv u'\pmod8$; this gives $2\cdot4=8$ classes, represented by $\pm1,\pm5,\pm2,\pm10$.
\end{solution}

\begin{solution}{ex:padicexp}
\idea Legendre's formula counts multiples of $p,p^2,\dots$; written in base $p$, each $\lfloor n/p^k\rfloor$ is $n$ with its last $k$ digits erased, and summing over $k$ telescopes into a formula with the digit sum. Then the exponential series converges if and only if its terms tend to $0$, which is a comparison of two linear rates.
\writeup For $x=0$ the exponential series equals $1$ and satisfies the stated radius condition. Assume $x\ne0$ for the valuation argument. Write $n=\sum_{i=0}^rc_ip^i$. Then $\lfloor n/p^k\rfloor=\sum_{i\ge k}c_ip^{i-k}$ and
\[
\begin{aligned}v_p(n!)&=\sum_{k\ge1}\sum_{i\ge k}c_ip^{i-k}=\sum_ic_i(p^{i-1}+\dots+1)\eqbr=\sum_ic_i\frac{p^i-1}{p-1}=\frac{n-s_p(n)}{p-1}.\end{aligned}
\]
So $v_p\left(\frac{x^n}{n!}\right)=n\left(v_p(x)-\frac1{p-1}\right)+\frac{s_p(n)}{p-1}$. If $v_p(x)>\frac1{p-1}$, this tends to $+\infty$ (as $s_p(n)\ge1$), so the terms tend to $0$ and the series converges. If $v_p(x)\le\frac1{p-1}$, take $n=p^m$ ($s_p(n)=1$): the valuation is $\le\frac1{p-1}$, so the terms do not tend to $0$. Since $v_p(x)>\frac1{p-1}$ means $|x|_p<p^{-1/(p-1)}$, we are done. In $\Q_2$: $|x|_2<\frac12$, that is, $x\in4\Z_2$.
\end{solution}

\begin{solution}{ex:hasse}
\idea As in Proposition~\ref{prop:noconic}: assume a primitive solution in $\Z_p$ (not all coordinates divisible by $p$) and look for a contradiction modulo $p$ or $p^2$. The good prime divides a coefficient and makes the form ``anisotropic'' modulo $p$: for $x^2-2y^2=3z^2$ try $p=3$, since $2$ is not a square modulo $3$; for $x^2+y^2=7z^2$ try $p=7$, since $-1$ is not a square modulo $7$.
\writeup For $x^2-2y^2=3z^2$ with a primitive solution in $\Z_3$: modulo $3$, $x^2\equiv2y^2$. If $3\nmid y$ then $2\equiv(x/y)^2$ is a square modulo $3$, false; so $3\mid y$, then $3\mid x$, then $9\mid3z^2$ and $3\mid z$: contradiction. For $x^2+y^2=7z^2$ in $\Z_7$: if $7\nmid y$ then $-1\equiv(x/y)^2\pmod7$, false since $7\equiv3\pmod4$; so $7\mid x,y$, then $49\mid7z^2$ and $7\mid z$: contradiction. In both cases there is no non-zero solution in $\Q_p$ for $p=3$, resp.\ $p=7$, hence none in $\Q$; both equations have real solutions, so the obstruction is purely $p$-adic.
\end{solution}

\begin{solution}{ex:teich}
\idea The key lemma says that raising to the power $p$ improves a congruence by one power of $p$. Starting from Fermat, $a^p\equiv a\pmod p$, it gives $a^{p^{n+1}}\equiv a^{p^n}\pmod{p^{n+1}}$, so the sequence is Cauchy. Its limit is a $(p-1)$-th root of unity congruent to $a$, which by uniqueness in Corollary~\ref{cor:teich} must be $\omega(a)$.
\writeup If $x=y+tp^k$, then $x^p=y^p+py^{p-1}tp^k+\sum_{j\ge2}\binom pjy^{p-j}t^jp^{jk}$. The second term is divisible by $p^{k+1}$, and so are the others: for $2\le j\le p-1$, $p\mid\binom pj$ and $jk\ge k+1$; for $j=p$, $pk\ge k+1$. By Fermat $a^p\equiv a\pmod p$; applying the lemma $n$ times, $a^{p^{n+1}}\equiv a^{p^n}\pmod{p^{n+1}}$. So $|a^{p^{n+1}}-a^{p^n}|_p\le p^{-n-1}$, the sequence is Cauchy, and it has a limit $\omega$. Then $\omega\equiv a^{p^n}\equiv a\pmod p$ (Fermat again), and $\omega^p=\lim a^{p^{n+1}}=\omega$, so $\omega^{p-1}=1$ since $\omega$ is a unit. By the uniqueness in Corollary~\ref{cor:teich}, $\omega=\omega(a)$. For $p=5$, $a=2$, modulo $125$: $2^{25}\bmod125$. We have $2^7=128\equiv3$, $2^{21}\equiv27$, $2^{25}=2^{21}\cdot16\equiv432\equiv57\pmod{125}$. So $\omega(2)\equiv57\pmod{125}$; check: $57^2=3249\equiv-1\pmod{125}$ ($3250=26\cdot125$), as it should be, since $\omega(2)$ is a primitive fourth root of unity ($2$ has order $4$ modulo $5$).
\end{solution}

\begin{solution}{ex:t1-random-valuation}
\idea Count survivors of each divisibility test, then subtract successive
survival probabilities and sum them for the expectation.
\writeup Of the $p^N$ equally likely residues, $p^{N-k}$ are divisible
by $p^k$, including zero. Thus $\operatorname{Pr}(V_N\ge k)=p^{-k}$.
For $0\le k<N$ the exact mass is
$p^{-k}-p^{-k-1}=(p-1)/p^{k+1}$; at $N$ it is $p^{-N}$,
the mass of zero.
Every value $v\in\{0,\dots,N\}$ satisfies
$v=\sum_{k=1}^N\mathbf1_{\{v\ge k\}}$, so finite linearity gives
\[
\mathbb E V_N=\sum_{k=1}^Np^{-k}=\frac{1-p^{-N}}{p-1}.
\]
For each fixed $k\ge0$, the mass eventually equals $(p-1)/p^{k+1}$.
These limiting masses sum to $1$, a geometric distribution on valuation values.
The sample spaces were residue rings. No equal mass on all positive integers
has been asserted; increasing precision and enlarging an interval are
different sampling models even when a particular limit agrees.
\end{solution}

\begin{solution}{ex:t1-minusone}
\idea A geometric series with ratio $5$ and first term $4$ sums to $4(5^n-1)/4=5^n-1$.
\writeup The sum of the first $n$ terms is $4\sum_{k=0}^{n-1}5^k=5^n-1\equiv -1\pmod{5^n}$. For $n=1,2,3$ the sums are $4$, $24$ and $124$. Therefore the coherent sequence of partial sums is the $5$-adic integer $-1$.
\end{solution}

\section*{Sources and further reading}

\begin{itemize}
\item F.~Q.~Gouv\^ea, \emph{$p$-adic Numbers: An Introduction}, 3rd ed., Universitext, Springer, 2020. The friendliest full introduction; our proof of Ostrowski's theorem follows the standard one given there.
\item N.~Koblitz, \emph{$p$-adic Numbers, $p$-adic Analysis, and Zeta-Functions}, 2nd ed., Graduate Texts in Mathematics~58, Springer, 1984. From the basics to $p$-adic zeta functions and Kummer's congruences.
\item J.-P.~Serre, \emph{A Course in Arithmetic}, Graduate Texts in Mathematics~7, Springer, 1973 (French original: \emph{Cours d'arithm\'etique}, PUF, 1970). Chapters~II--IV: $\Q_p$, the Hilbert symbol and Hasse--Minkowski, in a hundred perfect pages.
\item S.~Katok, \emph{$p$-adic Analysis Compared with Real}, Student Mathematical Library~37, AMS, 2007. Real and $p$-adic analysis side by side.
\item J.~W.~S.~Cassels, \emph{Local Fields}, London Mathematical Society Student Texts~3, Cambridge University Press, 1986. Includes Skolem--Mahler--Lech and many applications.
\item E.~S.~Selmer, ``The Diophantine equation $ax^3+by^3+cz^3=0$'', \emph{Acta Mathematica} 85 (1951), 203--362. The counterexample to the Hasse principle for cubics.
\item K.~Hensel, \emph{\"Uber eine neue Begr\"undung der Theorie der algebraischen Zahlen}, Jahresber.\ Deutsch.\ Math.-Verein.\ 6 (1897), 83--88.
\end{itemize}
